%
% EBR-III -- op:ebr3-assemble / ebr3-c full draft (local artifact).
% Portable amsart (venue class swapped per ebr3-f runbook; etna.cls/J.Algebra
% style applied only at the operator's chosen venue). Compile: pdflatex twice
% (inline thebibliography, no bibtex/latexmk needed on host).
% NOTHING in this file is deposited, committed, or submitted.
%
\documentclass[11pt,reqno]{amsart}

\usepackage{amsmath}
\usepackage{amssymb}
\usepackage{amsthm}
\usepackage{booktabs}
\usepackage{array}
\usepackage[hidelinks]{hyperref}
\usepackage{xurl}
\usepackage{microtype}

\setlength{\emergencystretch}{2em}

% Reproducible-PDF settings: omit embedded timestamps and fix the trailer /ID
% so the PDF byte-hash is stable across builds (used by the repro manifest;
% combine with SOURCE_DATE_EPOCH at build time).
\ifdefined\pdfinfoomitdate\pdfinfoomitdate=1\fi
\ifdefined\pdftrailerid\pdftrailerid{}\fi
\ifdefined\pdfsuppressptexinfo\pdfsuppressptexinfo=-1\fi

\theoremstyle{plain}
\newtheorem{theorem}{Theorem}
\newtheorem{proposition}[theorem]{Proposition}
\newtheorem{lemma}[theorem]{Lemma}
\newtheorem{corollary}[theorem]{Corollary}
\theoremstyle{definition}
\newtheorem{conjecture}[theorem]{Conjecture}
\theoremstyle{remark}
\newtheorem{remark}[theorem]{Remark}

\newcommand{\Q}{\mathbb{Q}}
\newcommand{\R}{\mathbb{R}}
\newcommand{\N}{\mathbb{N}}
\newcommand{\C}{\mathbb{C}}
\newcommand{\Z}{\mathbb{Z}}
\newcommand{\Fp}{\mathbb{F}_{p}}
\newcommand{\GGal}{G_{\mathrm{Gal}}}
\newcommand{\SL}{\operatorname{SL}}
\newcommand{\GL}{\operatorname{GL}}
\newcommand{\Sp}{\operatorname{Sp}}
\newcommand{\SO}{\operatorname{SO}}
\newcommand{\Sym}{\operatorname{Sym}}
\newcommand{\Hom}{\operatorname{Hom}}
\newcommand{\End}{\operatorname{End}}
\DeclareMathOperator{\Det}{det}
\newcommand{\Ltwo}{L_{2}}
\newcommand{\Cebr}{C_{\mathrm{EBR}}}
\newcommand{\claim}[1]{\textsf{\small[#1]}}

\begin{document}

\title[The EBR operator at $d=2$ has Galois group $\SL_4$]{The EBR operator
at $d=2$ is not a $G$-operator and its differential Galois group is
$\SL(4)$: arithmetic, irreducibility, and the non-Liouvillian corollary}

\author{Papanokechi}
\thanks{Independent Researcher (ORCID 0009-0000-6192-8273). This is a
preprint; nothing herein is a claim of priority. Every load-bearing statement
is graded \textsc{proven}/\textsc{structural}/\textsc{verified}/%
\textsc{conjectured} (\S\ref{sec:grades}); \textsc{proven} is reserved for
the Lean~4 cores of \S\ref{sec:lean}.}

\date{\today}

\subjclass[2020]{34M35, 12H05, 34M03, 33C20, 68V20, 11J81}

\keywords{differential Galois theory, connection coefficient, irregular
singularity, $G$-function, $p$-curvature, Aschbacher classes, Liouvillian
solutions, integer relation, PSLQ, Lean, Mathlib}

\begin{abstract}
For the positive-coefficient polynomial continued fraction with
$b(n)=3n^{2}+n+1$, the Edge--Borel object $G(s)=\sum_{n\ge0}Q_{n}s^{n}/(2n)!$
is annihilated by an explicit order-$4$ linear differential operator $\Ltwo$
over $\C(s)$ with regular singular points at $s=0$ and $s=R=4/3$ and an
\emph{irregular} singularity at $s=\infty$ of slope $1/4$. We establish four
things about $\Ltwo$ at increasing evidential strength. First, in the Borel
normalization $g_{n}=Q_{n}/(2n)!$ the operator is \emph{not} the minimal
operator of a $G$-function: the Taylor denominators retain essentially the
full factorial, the $p$-curvature is non-nilpotent for every tested prime,
and the irregular point at infinity is incompatible with the
Andr\'e--Chudnovsky--Katz trichotomy. Second, $\Ltwo$ is irreducible and
minimal, by two independent arguments. Third, its local monodromy at $R$ is a
\emph{semisimple} complex pseudo-reflection with spectrum
$\{1,1,1,e^{i\pi/3}\}$ --- correcting an earlier ``resonance logarithm''
reading, consistent with the criterion that a $-\gamma$ logarithm requires
$\gamma\in\Z$ whereas here $\gamma=11/6$. Fourth, and centrally, the identity
component of the differential Galois group is
$\GGal(\Ltwo)^{\circ}=\SL_{4}$; we obtain this by eliminating every class of
the maximal-subgroup (Aschbacher) classification of algebraic subgroups of
$\SL_{4}$, with the imprimitive class closed unconditionally through a unique
quadratic-character analysis and an \emph{a priori} degree bound that makes
the underlying rational-solution searches complete. The immediate corollary is
that $\Ltwo$ has no Liouvillian solutions. The connection coefficient is
computed by a third independent channel to $169$ stable digits; a
positive-control-validated integer-relation battery finds, to that precision,
that it is neither an elementary $\Gamma$-quotient nor a low-height
combination of standard constants nor algebraic of small degree. We state with
care what these results do \emph{not} give: a large Galois group does not
imply the connection coefficient is transcendental, connection coefficients
are not differential-Galois invariants, and the transcendence question
therefore remains a conjecture whose only route is a separate period analysis.
The degree-bound arithmetic, the pullback exponent arithmetic, the
eigenvalue-parity bookkeeping, and the parity of the relevant $4$-cycle are
formalized in Lean~4/Mathlib with audited axiom cones.
\end{abstract}

\maketitle

\section{Introduction}\label{sec:intro}

\subsection{Background and the object of study}
This paper continues the Edge--Borel Radius (EBR) program. EBR-I established
that for a positive-coefficient polynomial continued fraction
$Q_{n}=b(n)Q_{n-1}+Q_{n-2}$ with $b(n)=\beta_{d}n^{d}+\dots$, the physical
Borel transform $G(s)=\sum_{n\ge0}Q_{n}s^{n}/(dn)!$ is holonomic with a
dominant singularity at $s=R=d^{d}/\beta_{d}$ and local form
$G(s)\sim A\,(1-s/R)^{-\gamma}$, with $\gamma=(d+1)/2+b_{d-1}/\beta_{d}$
\cite{EBRI}. EBR-II identified the amplitude $A$ as a \emph{connection
coefficient} of the order-$2d$ operator annihilating $G$ --- the entry of the
connection matrix linking the exponent-$0$ local solution at the origin to the
dominant exponent-$(-\gamma)$ solution at $s=R$ --- and conjectured that this
coefficient is generically transcendental, with a ``rigidity dividing line''
separating the rigid case $d=1$ from the non-rigid cases $d\ge2$
\cite{EBRII}.

We fix the smallest non-rigid case. Throughout, $d=2$ and
\begin{equation}\label{eq:family}
  b(n)=3n^{2}+n+1,\qquad \beta_{2}=3,\qquad
  g_{n}=\frac{Q_{n}}{(2n)!},\qquad G(s)=\sum_{n\ge0}g_{n}s^{n},
\end{equation}
so that $R=4/3$ and $\gamma=11/6$. The annihilating operator, written with
$\theta=s\,d/ds$ as $L=P_{0}(\theta)-s\,P_{1}(\theta+1)-s^{2}$ and cleared to
the form $\sum_{k=0}^{4}a_{k}(s)\,d^{k}/ds^{k}$, has
\begin{equation}\label{eq:operator}
\begin{gathered}
  a_{4}=4s^{4}(4-3s),\quad a_{3}=48s^{3}-94s^{4},\quad a_{2}=12s^{2}-156s^{3},\\
  a_{1}=-30s^{2},\quad a_{0}=-s^{2}.
\end{gathered}
\end{equation}
We call this operator $\Ltwo$. Its leading coefficient
$a_{4}=d^{d}s^{2d}(d^{d}-\beta_{d}s)$ matches EBR-I, and its finite singular
locus is exactly $\{0,R\}$ \claim{CC1-OP-L2}.

\subsection{Results and their evidential status}\label{sec:grades}
We organize the paper by a four-class evidence discipline, stated here so that
the reader can audit each assertion against the kind of evidence that supports
it. \textsc{proven} is reserved for results carrying a machine-checked Lean~4
proof with an audited axiom cone; \textsc{structural} denotes a complete
mathematical argument verified by hand and computer algebra but not
formalized; \textsc{verified} denotes a reproducible high-precision or exact
symbolic computation; \textsc{conjectured} denotes an open statement. A grade
table appears in \S\ref{sec:gradetable}. The principal results are:

\begin{itemize}
\item[(R1)] \emph{Arithmetic nature.} In the normalization \eqref{eq:family},
$\Ltwo$ is not the minimal operator of a $G$-function (\S\ref{sec:gfun};
\textsc{verified}/\textsc{structural}).
\item[(R2)] \emph{Irreducibility.} $\Ltwo$ is irreducible and is the minimal
annihilator of $G$ (\S\ref{sec:irred}; \textsc{structural}, two routes).
\item[(R3)] \emph{Local monodromy.} The monodromy at $R$ is a semisimple
pseudo-reflection with spectrum $\{1,1,1,e^{i\pi/3}\}$ (\S\ref{sec:local};
\textsc{verified}), correcting an earlier reading (\S\ref{sec:erratum}).
\item[(R4)] \emph{Main theorem.} $\GGal(\Ltwo)^{\circ}=\SL_{4}$
(\S\ref{sec:main}; \textsc{structural}, bound-complete).
\item[(R5)] \emph{Corollary.} $\Ltwo$ has no Liouvillian solutions
(\S\ref{sec:bridge}; \textsc{verified}-by-citation applied to (R4)).
\item[(R6)] \emph{Connection data and a null.} The connection coefficient is
computed to $169$ digits and is, to that precision, not an elementary
$\Gamma$-quotient, not a low-height constant combination, and not algebraic of
small degree (\S\ref{sec:numerics}; \textsc{verified}).
\end{itemize}

\subsection{What we do not claim}\label{sec:nonclaim}
We are explicit about the limits of the method, since the temptation to
overread a large Galois group is precisely the error the EBR program guards
against.
\begin{quote}\itshape
Non-rigidity ($P=d-1>0$) does not imply that the connection coefficient is
transcendental; a large differential Galois group does not imply that the
connection coefficient is transcendental. The present work targets the
\emph{group} only. Connection coefficients are not differential-Galois
invariants, and the transcendence of the connection coefficient remains a
conjecture whose only route is a separate period analysis.
\end{quote}
Accordingly the abstract and this introduction state the $\SL_{4}$ theorem and
the not-a-$G$-function result, and no transcendence claim of any kind. The
$G$-function statement is always qualified by its normalization scope; the
group statement is for $d=2$, with general $d$ flagged open
(\S\ref{sec:open}); the explicit numerical Stokes multipliers at the ramified
irregular point are out of scope, for the reason given in
\S\ref{sec:stokes}.

\section{The operator and its singular structure}\label{sec:operator}

\subsection{Riemann data at the regular points}
At $s=0$ and $s=R=4/3$ the operator $\Ltwo$ is regular singular. The
indicial computation gives
\begin{equation}\label{eq:riemann}
\begin{gathered}
  \exp_{0}=\Big\{\tfrac{j}{d}:j=0,\dots,2d-1\Big\}=\{0,\tfrac12,1,\tfrac32\},\\
  \exp_{R}=\{0,1,2,-\gamma\}=\{0,1,2,-\tfrac{11}{6}\},
\end{gathered}
\end{equation}
with $\gamma=(d+1)/2+b_{d-1}/\beta_{d}=11/6$ \claim{CC1-RIEMANN}. The dominant
exponent $-\gamma$ differs from each of the integers $\{0,1,2\}$ by a
non-integer, a fact that will decide the local monodromy type in
\S\ref{sec:local}.

\subsection{The irregular point at infinity}\label{sec:inf}
The point at infinity is irregular. The Fuchs degree bound
$\deg a_{k}\le \deg a_{4}-(4-k)$ holds for $k=1,\dots,4$ but \emph{fails} at
$k=0$, where $\deg a_{0}=2>1$; the violation is caused exactly by the term
$a_{0}=-s^{2}$. The Newton polygon at infinity has a single edge from $(0,2)$
to $(4,5)$ of slope $3/4$, with exact edge polynomial $-12c^{4}-1$; its roots
$c^{4}=-1/12$ have $|c|=(1/12)^{1/4}=0.53728\ldots$ and arguments an odd
multiple of $\pi/4$. There is a single Puiseux cycle of length $4$, so the
ramification is exactly $4$, the slope is $1/4$, the integer Poincar\'e rank is
$1$, and the four determining factors
\begin{equation}\label{eq:detfactors}
  E_{k}(s)=\exp\!\big(4c\,\zeta_{4}^{k}\,s^{1/4}\big),\qquad
  c^{4}=-\tfrac{1}{12},\quad \zeta_{4}=i,\quad k=0,1,2,3,
\end{equation}
form one transitive orbit under the cyclic order-$4$ ramification group. The
prediction is corroborated numerically: at $s=10^{14}$ the four scaled
symbol-roots $w\,s^{3/4}$ have modulus tending to $0.53728$ and arguments
tending to $\{\pm\pi/4,\pm3\pi/4\}$ \claim{CC1-INF-IRREGULAR}. This corrects
the earlier ``three regular singular points'' reading of the operator.

\subsection{Accessory count and non-rigidity}
Because infinity is irregular, the Fuchsian accessory count
$N_{\mathrm{acc}}=(2d-1)(d-1)$ of EBR-II is inapplicable; the corrected
index-of-rigidity count is $P=d-1$, which at $d=2$ equals $1$. We do not
re-derive the full Katz rigidity index with the irregular point here; this is
flagged as a residual gap (\S\ref{sec:open}). The point that matters
downstream is robust to the discrepancy: the count is positive at $d=2$ on
either reading, so the connection problem is non-rigid \claim{CC1-ACCESSORY}.

\section{Arithmetic nature: not a $G$-operator}\label{sec:gfun}

We emphasize at the outset the \emph{normalization scope} of this section: all
statements are about the Borel normalization $g_{n}=Q_{n}/(2n)!$ of
\eqref{eq:family} and do not extend to arbitrary rescalings of the coefficient
sequence (multiplying by $(2n)!$, for instance, returns the integer series
$\sum Q_{n}s^{n}$ of radius $0$, a different object). With that scope fixed:

\subsection{Integrality and denominator growth}
The numerators are integers: $Q_{0}=1$, $Q_{1}=5$, and
$Q_{n}=(3n^{2}+n+1)Q_{n-1}+Q_{n-2}$ has integer coefficients, so $Q_{n}\in\Z$
(checked to $n\le2000$), whence $\operatorname{denom}(g_{n})\mid(2n)!$ and
$q_{N}:=\operatorname{lcm}_{n\le N}\operatorname{denom}(g_{n})\mid(2N)!$
\claim{CC2-0-QINT}. The content of the section is that this divisibility is
\emph{not} far from an equality: $q_{N}$ grows superlinearly, the fit
$q_{N}\sim N\log N$ giving $R^{2}=0.9999983$ against $R^{2}=0.998858$ for a
linear law, and $\log_{10}q_{N}/\log_{10}(2N)!\approx1.0000$. The Taylor
denominators retain essentially the entire factorial; there is no geometric
denominator collapse, so $G$ is not a $G$-function in the
Andr\'e--Bombieri sense, but a Gevrey/Borel object \claim{CC2-0-GFUNC}.

\subsection{Two confirming channels}
The $p$-curvature $\psi_{p}$ of $\Ltwo$ over $\Fp(s)$ is non-nilpotent for
every tested prime $p\in\{5,7,11,13,17,19,23,29\}$, at three generic
specializations each (the primes $2,3$ are excluded as degenerate for the
leading coefficient). Since global nilpotence is equivalent to nilpotence of
$\psi_{p}$ for almost all $p$, the operator is not globally nilpotent, which by
the Chudnovsky--Andr\'e theorem again excludes the $G$-function property and
matches the irregular-infinity prediction of Katz \claim{CC2-0-PCURV}. Indeed
the Andr\'e--Chudnovsky--Katz results yield a clean trichotomy: of the three
statements ``$G$ is a $G$-function'', ``$\Ltwo$ is minimal/irreducible'', and
``$\infty$ is irregular'', at most two can hold. We establish the latter two
(\S\ref{sec:irred} and \S\ref{sec:inf}), so the first fails
\claim{CC2-0-TRICHOTOMY}.

\subsection{Consequence for the period route}
A direct corollary concerns the transcendence question. Because $G$ itself is
not a $G$-function, no unconditional transcendence statement about the
connection coefficient can be drawn from ``$G$ is a $G$-function with the right
growth''. The $G$-function/period machinery may still bear on the coefficient
as a connection datum of the order-$4$ operator that is Fuchsian at $\{0,R\}$
and irregular at infinity, but establishing that is the burden of the separate
period analysis and cannot be inherited as assumed \claim{CC2-0-CC3-AMEND}.

\section{Irreducibility and minimality}\label{sec:irred}

The falsification target ``$\Ltwo$ is reducible over $\C(s)$'' is refuted by
two independent arguments.

\subsection{Formal module at infinity}
At the irregular point the formal $\C((s^{-1}))$-module of $\Ltwo$ has a single
slope $1/4$ whose four determining factors \eqref{eq:detfactors} form one
transitive orbit under the cyclic order-$4$ ramification group. A proper right
factor over $\C(s)$ would have a $\C(s)$-rational, hence ramification-stable,
set of determining factors; but the only subsets of a transitive size-$4$
orbit stable under the generating $4$-cycle are the empty set and the whole
set. Hence $\Ltwo$ has no proper right factor and is the minimal annihilator of
$G$ (order $4=2d$). The local-to-global inference uses Turrittin's formal
classification and the formal-module description of
\cite[Ch.~3]{vdPutSinger}, and is graded \textsc{structural} (computer algebra
and hand argument; an independent \texttt{DFactor}-style confirmation was
unavailable on the host) \claim{CC1-IRREDUCIBLE}.

\subsection{Exhaustive low-order factor exclusion}
Independently, slope additivity at infinity (the Newton polygon of a product is
the concatenation of the factors' polygons) reduces the existence of an
order-$m$ right factor to the existence of a $\sigma$-stable size-$m$ subset of
the determining-factor orbit, where $\sigma=(0\,1\,2\,3)$ generates the
ramification group. A finite enumeration shows the stable subset sizes are
$\{0,4\}$, so no order-$1$, $2$, or $3$ right factor exists; for order $1$,
each factor in \eqref{eq:detfactors} is ramified, so there is no single-valued
$\exp\!\int r$ with $r\in\C(s)$. This recovers irreducibility by an
independent route \claim{CC2-0-FACTOR}. The two routes are logically
independent --- one local-analytic at infinity, one a finite combinatorial
enumeration of $\sigma$-stable orbits --- and neither relies on a
decision-procedure factorization; we grade irreducibility \textsc{structural}
on the strength of their agreement, and it is this irreducibility that the
degree-bound completeness of \S\ref{sec:bounds} later makes load-bearing.

\section{Local monodromy and the semisimple correction}\label{sec:local}

\subsection{The monodromy at $R$ is a semisimple pseudo-reflection}
The local monodromy at $R$ has eigenvalues $e^{2\pi i\rho}$ for
$\rho\in\exp_{R}=\{0,1,2,-\tfrac{11}{6}\}$, that is $\{1,1,1,e^{i\pi/3}\}$
(using $e^{-2\pi i\cdot11/6}=e^{i\pi/3}$). Although the eigenvalue $1$ occurs
with multiplicity $3$ and there is an integer-exponent resonance among
$\{0,1,2\}$, no logarithm forms: an exact symbolic Frobenius solve at $R$ shows
that all resonance obstructions on the eigenvalue-$1$ tower vanish, giving
three independent logarithm-free solutions, while an independent numerical
computation gives $\operatorname{rank}(M_{R}-I)=1$ (one singular value
$2.159$, the rest below $10^{-30}$). Thus $M_{R}$ is a semisimple complex
pseudo-reflection of order $6$ with spectrum $\{1,1,1,e^{i\pi/3}\}$, and
$M_{0}$ is semisimple with spectrum $\{1,-1,1,-1\}$
\claim{CC2-2D-JORDAN}\claim{CC4-1-MONODROMY}. The exact-arithmetic resonance
residuals were subsequently pushed to $169$ digits, confirming the absence of
a logarithm to that precision \claim{CC4-1-NOLOG}.

\subsection{Erratum}\label{sec:erratum}
This corrects an earlier narration that spoke of a ``unipotent part'' or
``resonance logarithm'' at $R$. The corrected statement is that $M_{R}$ is a
semisimple pseudo-reflection with no logarithm, consistent with the criterion
of \cite{EBRII} that a $-\gamma$ resonance logarithm requires $\gamma\in\Z$;
here $\gamma=11/6\notin\Z$. The superseded reading appeared in the narration of
EBR-II (concept DOI \texttt{10.5281/zenodo.20566465}, current version
\texttt{10.5281/zenodo.20571232}, v1.2); the load-bearing claims of that work
--- the exponents, the irreducibility, and the irregularity at infinity --- are
unaffected, only the adjective on $M_{R}$ \claim{CC4-ERR-1}.

\section{The differential Galois group}\label{sec:main}

We compute $\GGal(\Ltwo)$ over $\C(s)$. The strategy is the
maximal-subgroup classification of algebraic subgroups of $\SL_{4}$ in the
differential-Galois setting, following the formulation of
\cite[Ch.~4]{vdPutSinger}; we use the Aschbacher class labels
$C_{1},\dots,C_{8},S$ as organizing vocabulary for the maximal closed subgroups
of $\SL_{4}$, and we stress that this is the algebraic-group classification,
not an application of Aschbacher's finite-group theorem. In the geometric
(Kleidman--Liebeck) numbering the classical-form subgroups $\SO_{4},\Sp_{4}$
are class $C_{8}$, the subfield class is $C_{5}$, the extraspecial-normalizer
class is $C_{6}$, and the tensor-induced class is $C_{7}$. Over the
algebraically closed field of constants in characteristic zero the subfield
class $C_{5}$ is empty; the tensor-induced class $C_{7}$ for the $2\otimes2$
decomposition shares its identity component $\SL_{2}\otimes\SL_{2}$ with $C_{4}$
and is excluded with it; and the extraspecial-type class $C_{6}$ is finite,
hence killed together with the other finite primitive classes by the
exponential torus. The operative algebraic classes are therefore
$C_{1},C_{2},C_{3},C_{4},C_{8},S$ together with the finite classes.

\subsection{Determinant and the $\SL_{4}$ normalization}
The Wronskian satisfies $w'/w=-a_{3}/a_{4}=(24-47s)/(2s(3s-4))
=-3/s+(29/2)/(4-3s)$, so $W=s^{-3}(4-3s)^{-29/6}$, with determinant exponents
$-3$ at $0$ and $-29/6$ at $R$ (each cross-checked by
$W_{\exp}=\sum\exp-\binom{4}{2}$ with residual $0$). The $\SL_{4}$-normalizing
twist $\chi=\Det^{1/4}=s^{-3/4}(4-3s)^{-29/24}$ is algebraic over $\C(s)$ (a
product of rational powers; not rational, not properly Liouvillian). After
twisting by $\chi$ the regular-point exponent multisets become
$\{3/4,5/4,7/4,9/4\}$ at $0$ and $\{29/24,53/24,77/24,-5/8\}$ at $R$, each
summing to $6$ \claim{CC2-1-TWIST}. The four determining-factor leading units
$c\cdot\{1,i,-1,-i\}$ at infinity generate a rank-$2$ lattice
$\Z\cdot1+\Z\cdot i$, so the formal exponential torus has dimension $2$; in
particular $\GGal(\Ltwo)^{\circ}$ contains a $2$-dimensional torus and is
infinite \claim{CC2-1-EXPTORUS}.

\subsection{No invariant form: $C_{8}$ ($\SO_{4}$, $\Sp_{4}$)}
\label{sec:selfdual}
A nondegenerate $G$-invariant bilinear form, even up to a rank-$1$ twist
$V\cong V^{*}\otimes\kappa$, would force the local eigenvalue multiset at $R$ to
equal $\kappa\cdot\{\text{inverses}\}$. The multiplicity-$3$ eigenvalue $1$
pins $\kappa=1$, and then $e^{i\pi/3}\ne e^{-i\pi/3}$. Hence the multiset
$\{1,1,1,e^{i\pi/3}\}$ is not inversion-closed up to any scalar, so $\Ltwo$
carries no invariant symmetric or alternating form, even up to twist: the group
lies in neither $\SO_{4}$ nor $\Sp_{4}$. This argument is exact and uses only
the $-11/6$ exponent \claim{CC2-2A-SELFDUAL}. It is confirmed constructively:
the exterior square $\Lambda^{2}\Ltwo$ has order $6$ (no symplectic
order-drop), splitting at infinity into a rank-$2$ slope-$0$ piece --- the
cancelling pairs $(0,2)$ and $(1,3)$, since $i^{0}+i^{2}=0$ and
$i^{1}+i^{3}=0$ --- and a rank-$4$ slope-$1/4$ remainder; the search for a
rational antisymmetric solution of $P'=-A^{\!\top}P-PA$ returns nullspace
dimension $0$ in two ansatz boxes (denominator $s^{a}(4-3s)^{b}$, numerator
degree $d$; $(a,b,d)=(4,4,10)$ and $(3,3,8)$), so the group is not in
$\Sp_{4}$ \claim{CC2-2B-SP4}. The same search for a symmetric solution returns
nullspace dimension $0$ in the same boxes, so the group is not in $\SO_{4}$
\claim{CC2-2C-SO4}.

\subsection{Tensor and tensor-induced classes: $C_{4}$, $C_{7}$, $S$}
A symmetric cube $\Sym^{3}W$ of a rank-$2$ module is self-dual up to twist,
contradicting \S\ref{sec:selfdual}; moreover a rank-$2$ module has ramification
at most $2$ and cannot carry slope $1/4$, and $\Sym^{3}$ would produce factors
$\{3\phi,\phi,-\phi,-3\phi\}$ of two distinct magnitudes rather than the single
equal-magnitude transitive $\Z/4$ orbit of $\Ltwo$. Hence the $\Sym^{3}\SL_{2}$
class $S$ is excluded. A tensor product $V_{1}\otimes V_{2}$ of two rank-$2$
modules (class $C_{4}$) is likewise self-dual up to twist and again cannot
produce slope $1/4$, so it too is excluded; the tensor-induced class $C_{7}$
(the wreath $\SL_{2}\wr S_{2}$ acting on the $2\otimes2$ decomposition) has the
same identity component $\SL_{2}\otimes\SL_{2}$ and is excluded by the same
argument \claim{CC2-2D-TENSOR}.

\subsection{The imprimitive classes $C_{2},C_{3}$ and primitivity}
\label{sec:primitive}
The remaining obstruction is imprimitivity: the monomial/$2$-block class
$C_{2}$ and the field-extension class $C_{3}$. These are not separable from the
local data alone, because the connected exponential torus fixes every block of
any system of imprimitivity, forcing blocks to be spans of weight lines, and
both the $4$-cycle formal monodromy and the semisimple pseudo-reflection
$M_{R}$ preserve the $2$-block system $\{\{L_{1},L_{3}\},\{L_{2},L_{4}\}\}$
\claim{CC2-2D-IMPRIM}. We close the class globally.

A quadratic character $\eta$ realizing an imprimitive structure ramifies at an
even number of singular points. Since the $R$-spectrum $\{1,1,1,e^{i\pi/3}\}$
is not negation-closed (no $-1$), $\eta$ must be unramified at $R$; evenness
then forces ramification exactly at $\{0,\infty\}$, i.e.\ $\eta=\sqrt{s}$, which
passes the necessary local tests because the spectra at $0$ ($\{1,-1,1,-1\}$)
and at $\infty$ (units $\{1,i,-1,-i\}$) are negation-closed. Thus $\sqrt{s}$ is
the unique candidate character \claim{CC4-A1-ETA-UNIQUE}. Furthermore the
formal monodromy at infinity is the $4$-cycle $(q_{0}\,q_{1}\,q_{2}\,q_{3})$ on
the determining factors (since $s^{1/4}\mapsto i\,s^{1/4}$), an odd permutation;
any induced structure embeds the permutation action in a transitive subgroup of
$S_{4}$ containing this $4$-cycle, hence in $\{C_{4},D_{4},S_{4}\}$, each of
which has an index-$2$ subgroup, i.e.\ a quadratic character, which by
uniqueness is $\sqrt{s}$. This reduces the $2$-block, the $4$-line monomial,
and the degree-$2/4$ field-extension classes all to the single $\sqrt{s}$ test
\claim{CC4-A1B-MONOMIAL-CLOSURE}.

That test is negative, by two independent searches. Over $\C(s)$, the search
for a rational intertwiner $\Phi$ realizing $M\cong M\otimes\sqrt{s}$ (solving
$\Phi'=A\Phi-\Phi A+(1/2s)\Phi$) returns nullspace dimension $0$ across boxes
$s^{a}(4-3s)^{b}$ of degree up to $20$ ($336$ unknowns at the largest); as a
calibration, the untwisted search returns $\End_{\C(s)}(M)=1$ (scalars only),
confirming absolute irreducibility and that the pipeline detects solutions when
they exist \claim{CC4-0-ROUTE2-PRIMITIVE}. Independently, the pullback under
$s=t^{2}$ gives $\tilde A(t)=2t\,A(t^{2})$, with $t=0$ exponents doubling to
$\{0,1,2,3\}$, slope $1/2$ at infinity, and two unramified points $\pm t_{R}$
swapped by the deck map; the eigenring of the pullback is $1$-dimensional
(scalars), so the pullback is irreducible over $\C(t)$ and $M$ is not induced
from the quadratic cover \claim{CC4-0-ROUTE1-PRIMITIVE}. The two routes agree
on primitivity.

\subsection{Completeness of the searches}\label{sec:bounds}
The two rational-solution searches are made complete by an \emph{a priori}
degree bound. Because $M_{0}$ and $M_{R}$ are semisimple (no logarithm), a
rational horizontal section is single-valued and every non-integer
$\End$-exponent branch is killed, so the pole order at a finite singular point
$p$ is at most $-\min$ of the integer $\End$-exponents at $p$ (these exponents
being pairwise differences of the $M$-exponents, shifted by the $\eta$-twist),
and the growth at infinity is bounded by the minimal integer slope-$0$ exponent
of the ramified formal monodromy. Computed exactly, the numerator-degree bound
for Route~2 is $B_{2}=1+2+0=3$ (pole $\le1$ at $0$, $\le2$ at $R$, no growth at
infinity), and for Route~1 it is $B_{1}=3+2+2=7$ (pole $\le3$ at $t=0$, $\le2$
at each $\pm t_{R}$). Both are at most $20$, and the boxes actually run ---
$[s(4-3s)^{2},\deg\le8]$ and $[t^{3}(3t^{2}-4)^{2},\deg\le10]$ --- strictly
contain $B_{2}$ and $B_{1}$. Hence the nulls of \S\ref{sec:primitive} are
bound-complete, not merely empirical \claim{CC4-0B-BOUNDS}.

\subsection{The classification and the main theorem}
\label{sec:asch}
We collect the eliminations.

\begin{center}
\small
\begin{tabular}{@{}ll>{\raggedright\arraybackslash}p{0.42\textwidth}l@{}}
\toprule
Class & Description & Eliminated by & Grade \\
\midrule
$C_{1}$ & reducible & \S\ref{sec:irred} (\claim{CC1-IRREDUCIBLE},
  \claim{CC2-0-FACTOR}) & S \\
$C_{6}$/fin. & extraspecial-type / finite & \S\ref{sec:main} torus (\claim{CC2-1-EXPTORUS}) & S \\
$C_{8}$ & orthogonal $\SO_{4}$ & \S\ref{sec:selfdual} (\claim{CC2-2C-SO4},
  \claim{CC2-2A-SELFDUAL}) & V/S \\
$C_{8}$ & symplectic $\Sp_{4}$ & \S\ref{sec:selfdual} (\claim{CC2-2B-SP4},
  \claim{CC2-2A-SELFDUAL}) & V/S \\
$S$ & $\Sym^{3}\SL_{2}$ & tensor/slope (\claim{CC2-2D-TENSOR}) & S \\
$C_{4},C_{7}$ & $\SL_{2}\otimes\SL_{2}$ & tensor/slope (\claim{CC2-2D-TENSOR}) & S \\
$C_{2}$ & monomial / $2$-block & $\sqrt{s}$ test (\claim{CC4-0-ROUTE2-PRIMITIVE},
  \claim{CC4-0-ROUTE1-PRIMITIVE}) & S \\
$C_{3}$ & field extension & $4$-cycle closure (\claim{CC4-A1B-MONOMIAL-CLOSURE})
  & S \\
$C_{5}$ & subfield & vacuous over $\overline{\C}$, char $0$ & S \\
\bottomrule
\end{tabular}
\end{center}

\begin{theorem}\label{thm:main}
For the operator $\Ltwo$ of \eqref{eq:operator}, the identity component of the
differential Galois group over $\C(s)$ is
\[
  \GGal(\Ltwo)^{\circ}=\SL_{4}.
\]
\end{theorem}

\begin{proof}[Evidential status of the proof]
The two imprimitivity-excluding routes ($\Hom(M,M\otimes\sqrt{s})=0$ over
$\C(s)$; eigenring $=$ scalars on the $s=t^{2}$ cover) agree and are
bound-complete by \S\ref{sec:bounds}; the unique-character analysis
(\claim{CC4-A1-ETA-UNIQUE}) and the $4$-cycle closure
(\claim{CC4-A1B-MONOMIAL-CLOSURE}) reduce every imprimitive class to the single
$\sqrt{s}$ test; the remaining maximal classes are excluded as tabulated. With
$C_{1},\dots,C_{8},S$ exhausted and the torus forcing the group to be infinite
and reductive, the identity component is $\SL_{4}$. The graded inputs are
\textsc{structural} (the routes, the character analysis, the bound) and
\textsc{verified}-by-citation (the classification framework of
\cite{vdPutSinger}); the theorem is therefore \textsc{structural} and, with the
bound of \S\ref{sec:bounds}, unconditional \claim{CC4-0-SL4-VERDICT}.
\end{proof}

\section{The non-Liouvillian corollary and the bridge}\label{sec:bridge}

\begin{corollary}\label{cor:liouville}
$\Ltwo$ has no Liouvillian solutions.
\end{corollary}

\begin{proof}
A linear operator has a Liouvillian solution if and only if the identity
component of its differential Galois group is solvable
\cite[Ch.~1,4]{vdPutSinger}. By Theorem~\ref{thm:main} the identity component
is $\SL_{4}$, which is connected, semisimple, and non-solvable; hence no
Liouvillian solution exists. This is a citation applied to (R4) and is graded
\textsc{verified}-by-citation $+$ \textsc{structural} application
\claim{CC4-2-BRIDGE}.
\end{proof}

We restate the discipline of \S\ref{sec:nonclaim} as the bridge to the open
problem. The group computation licenses exactly the non-Liouvillian corollary
and nothing about the arithmetic nature of the connection coefficient.
Connection coefficients are not differential-Galois invariants; the
transcendence of the connection coefficient remains a conjecture, and its only
route is the separate period analysis, which must be argued afresh because the
$G$-function clause is retired by \S\ref{sec:gfun}.

\section{Connection data and the integer-relation null}\label{sec:numerics}

\subsection{The connection coefficient to $169$ digits}
We computed the connection coefficient by a third independent channel --- a
full $4\times4$ local fundamental system at $R$ in the basis
$\{u_{-11/6},u_{0},u_{1},u_{2}\}$, matched to the continued series $G$ and its
first three derivatives at an interior point --- obtaining the amplitude
\[
  A=2.87436850995728075034795582039806005699192841843121\ldots
\]
and the coefficient-asymptotics prefactor
\[
  \Cebr=\frac{A}{\Gamma(11/6)}
       =3.055706807890481365701912201727\ldots,
\]
the latter appearing in $g_{n}\sim\Cebr\,R^{-n}n^{\gamma-1}$. The channel is
self-validated by multi-point agreement ($s^{*}=1.0$ versus $0.9$, agreeing to
$169$ digits) and multi-precision agreement (working precision $170$ versus
$210$, agreeing to $169$ digits); a Richardson channel on the coefficient
asymptotics independently confirms $77$ digits, the clean $1/n$ extrapolation
being itself a re-confirmation of the absence of a logarithmic tower at $R$. We
state honestly that these are \texttt{mpmath} floating values, not formal
interval enclosures; the grade is \textsc{verified}, not \textsc{proven}
\claim{CC4-1-C-120D}.

\subsection{Stokes structure, and what is out of scope}\label{sec:stokes}
At the irregular point the four determining factors $q_{k}=\gamma_{k}s^{1/4}$
with $\gamma_{k}^{4}=-1/12$ give eight singular directions in the
$w=s^{1/4}$ plane, spaced $\pi/4$, and a formal monodromy equal to a single
$4$-cycle (permutation sign $-1$) times $\operatorname{diag}(e^{2\pi i\mu_{k}})$.
Determinant consistency of the cyclic relation $M_{0}M_{R}M_{\infty}=I$ forces
$\Det M_{\infty}=e^{-i\pi/3}$, i.e.\ $\sum_{k}\mu_{k}\equiv1/3\pmod 1$, a
constraint any Stokes computation must satisfy. We are explicit that the
\emph{numerical} Stokes multipliers at the ramified slope $1/4$ require
Borel--Laplace multisummation that we did not complete on the host; we report
the layout and the determinant constraint and do not substitute a weaker proxy
for the multipliers \claim{CC4-1-STOKES}. This is the one computational leg of
the program left open here (\S\ref{sec:open}).

\subsection{The integer-relation null as exclusion bounds}\label{sec:null}
We tested the falsification target ``the connection coefficient is an
elementary $\Gamma$-quotient, a low-height constant combination, or
algebraic''. The targets $\Cebr$ and $A$ are known to $169$ digits; we ran a
PSLQ battery at working precision $169$ with acceptance tolerance $10^{-150}$,
leaving roughly $19$ guard digits. The battery's validity is demonstrated by
three positive controls, all detected with the expected integer vectors: the
$\Gamma$-reflection $\Gamma(1/3)\Gamma(2/3)=2\pi/\sqrt3$, the algebraic relation
$x^{2}-2=0$ for $\sqrt2$, and the linear relation $2\pi+3\log2$.

For a basis of $n$ terms, a spurious relation of integer height $H$ produces a
residual of order $H^{-(n-1)}$; for the search to be limited by the searched
height rather than by precision we require $H^{-(n-1)}$ to remain above the
tolerance, i.e.\ the precision-detection ceiling is
$H_{\max}\approx10^{150/(n-1)}$. In every tier below the searched height lies
well within this ceiling, so each null is a genuine exclusion bound: \emph{any}
integer relation must have height exceeding the searched $H$.

\begin{center}
\footnotesize
\setlength{\tabcolsep}{4pt}
\begin{tabular}{@{}l>{\raggedright\arraybackslash}p{0.32\textwidth}lll@{}}
\toprule
Tier & basis ($n$ terms) & searched $H$ & ceiling $10^{150/(n-1)}$ & verdict \\
\midrule
$\Gamma$-monomial & reflection-normalized $\Gamma$ at $(1/24)\Z$, indep.\ ($n{=}8$)
  & $10^{9}$ & $10^{21.4}$ & null \\
$\Gamma$-monomial & overcomplete $11$-$\Gamma$ basis ($n{=}15$)
  & $10^{9}$ & $10^{10.7}$ & inter-$\Gamma$ only \\
constants & $\{1,\pi,\log2,\log3,$ $\gamma_{E},\text{Cat},\zeta_3,\zeta_5\}$ ($n{=}9$)
  & $10^{12}$ & $10^{18.8}$ & null \\
algebraic & minimal polynomial $\deg\le8$ ($n\le9$)
  & $10^{6}$ & $10^{18.8}$ & null \\
\bottomrule
\end{tabular}
\end{center}

In the reflection-normalized independent $\Gamma$-basis there is no relation at
all; in the overcomplete $11$-$\Gamma$ basis PSLQ recovers only genuine Gauss
multiplication/reflection identities among the $\Gamma$-values, each with
\emph{coefficient zero on the target} (residuals $\sim10^{-170}$) --- an
internal control showing the battery sees real $\Gamma$-structure without
entangling the target. To $169$ digits, then, $\Cebr$ and $A$ are not
elementary $\Gamma$-quotients \claim{EBR3-B-GAMMA}, not integer-linear
combinations of $\{1,\pi,\log2,\log3,\gamma_{E},\text{Catalan},\zeta(3),
\zeta(5)\}$ \claim{EBR3-B-CONST}, and not algebraic of degree $\le8$ and height
$\le10^{6}$ \claim{EBR3-B-ALG}. This strengthens the $\ge33$-digit
$\Gamma$-quotient null of EBR-II to $169$ digits. It is evidence of
non-elementarity; it is not a proof of transcendence, and we do not present it
as one.

\section{Lean cores and the axiom audit}\label{sec:lean}

The finitary, arithmetic sub-claims that admit machine-checked proofs are
formalized in Lean~4 with Mathlib (\texttt{leanprover/lean4} v4.30.0; Mathlib
rev \texttt{c5ea0035}). Four cores are proved: the degree-bound arithmetic
$B_{2}=1+2+0=3$, $B_{1}=3+2+2=7$, both $\le20$ (\claim{CC4-LEAN-BOUNDS}); the
pullback exponent integrality, that $s=t^{2}$ doubles
$\{0,\tfrac12,1,\tfrac32\}$ to the integers $\{0,1,2,3\}$
(\claim{CC4-LEAN-PULLBACK}); the eigenvalue-parity bookkeeping, that the
$R$-spectrum is not negation-closed while the $0$-spectrum is, encoded by angles
in $\Z/6$ with negation $k\mapsto k+3$ (\claim{CC4-LEAN-PARITY}); and the parity
core, that the $4$-cycle \texttt{finRotate 4} is odd and so lies outside the
alternating group, the group-theoretic heart of the $4$-cycle closure
(\claim{CC4-LEAN-A1B}). A \verb|#print axioms| audit confirms that every
declaration depends only on a cone contained in
$\{\texttt{propext},\texttt{Classical.choice},\texttt{Quot.sound}\}$, with no
\texttt{sorryAx}; three of the bound lemmas depend on no axioms at all. These
are the only \textsc{proven} results in the paper.

\section{Four-class grade table}\label{sec:gradetable}

\begin{center}
\small
\begin{tabular}{@{}p{0.42\textwidth}l>{\raggedright\arraybackslash}p{0.34\textwidth}@{}}
\toprule
Result & Grade & Evidence \\
\midrule
$\Ltwo$ explicit, singular locus $\{0,R\}$ & V & \claim{CC1-OP-L2} \\
Riemann data \eqref{eq:riemann} & V & \claim{CC1-RIEMANN} \\
$\infty$ irregular, slope $1/4$ & V & \claim{CC1-INF-IRREGULAR} \\
Not a $G$-function (normalization-scoped) & V & \claim{CC2-0-GFUNC},
  \claim{CC2-0-QINT} \\
$p$-curvature non-nilpotent & V & \claim{CC2-0-PCURV} \\
Trichotomy excludes $G$-function & S & \claim{CC2-0-TRICHOTOMY} \\
Irreducible/minimal (two routes) & S & \claim{CC1-IRREDUCIBLE},
  \claim{CC2-0-FACTOR} \\
$M_{R}$ semisimple pseudo-reflection & V & \claim{CC2-2D-JORDAN},
  \claim{CC4-1-NOLOG} \\
No invariant form ($\not\le\SO_{4},\Sp_{4}$) & V/S & \claim{CC2-2A-SELFDUAL},
  \claim{CC2-2B-SP4}, \claim{CC2-2C-SO4} \\
Tensor classes excluded & S & \claim{CC2-2D-TENSOR} \\
$\eta=\sqrt{s}$ unique; $4$-cycle closure & S & \claim{CC4-A1-ETA-UNIQUE},
  \claim{CC4-A1B-MONOMIAL-CLOSURE} \\
Primitive (two agreeing routes) & S & \claim{CC4-0-ROUTE1-PRIMITIVE},
  \claim{CC4-0-ROUTE2-PRIMITIVE} \\
Degree bounds $B_2{=}3,B_1{=}7\le20$ & S & \claim{CC4-0B-BOUNDS} \\
\textbf{Main:} $\GGal(\Ltwo)^{\circ}=\SL_{4}$ & S & \claim{CC4-0-SL4-VERDICT} \\
No Liouvillian solutions & V/S & \claim{CC4-2-BRIDGE} \\
Connection coefficient to $169$ digits & V & \claim{CC4-1-C-120D} \\
$\Gamma$-quotient / constant / algebraic nulls & V & \claim{EBR3-B-GAMMA},
  \claim{EBR3-B-CONST}, \claim{EBR3-B-ALG} \\
Lean cores (4) & P & \claim{CC4-LEAN-BOUNDS}, \claim{CC4-LEAN-PULLBACK},
  \claim{CC4-LEAN-PARITY}, \claim{CC4-LEAN-A1B} \\
Transcendence of the connection coefficient & C & \claim{CC1-DISCIPLINE} \\
\bottomrule
\end{tabular}
\end{center}

\section{Open problems}\label{sec:open}

We list the open directions, each tagged with the program op-code under which
it is being pursued.
\begin{enumerate}
\item \emph{Transcendence of the connection coefficient} (\texttt{op:cc-3}).
The sole route is via a period-theoretic analysis, argued from first principles since
the $G$-function clause is retired (\S\ref{sec:gfun}). The $169$-digit null of
\S\ref{sec:null} is the current evidence of non-elementarity; the
$169$-digit value, the exact monodromy, and the Stokes layout are the inputs.
\item \emph{Numerical Stokes multipliers} at the ramified slope $1/4$, via
Borel--Laplace multisummation (\S\ref{sec:stokes}).
\item \emph{General $d$.} Only $d=2$ is settled here; the irreducibility and
the group computation at $d\ge3$ require the order-$2d$ analogues of every step.
Three ingredients are specific to $d=2$ and do not transfer for free: the
determining-factor orbit at infinity is a single $4$-cycle (a transitive
$\Z/2d$-orbit at general $d$, so the slope-additivity factor enumeration grows),
the target group becomes $\SL_{2d}$ with its own and larger maximal-subgroup
list, and the \emph{a priori} degree bounds of \S\ref{sec:bounds} must be
recomputed from the order-$2d$ exponent data. The arithmetic conclusions are
expected to persist but are not established here.
\item \emph{The rigidity index} with the irregular point, completing the
accessory count beyond the robust positivity used in
\S\ref{sec:operator}.
\item \emph{The rigidity dividing-line conjecture.} The precise statement ---
that the connection coefficient is elementary exactly on the rigid ($d=1$) or
resonant/reducible members and a genuine transcendental period generically for
$d\ge2$ --- remains conjectural; the present work supplies the group side at
$d=2$ but not the arithmetic side.
\end{enumerate}

\section*{Reproducibility}
All computations are reproducible from the accompanying package: the symbolic
operator and singular-structure scripts, the $p$-curvature and factor-exclusion
scripts, the twist and invariant-form searches, the primitivity routes and
degree bounds, the $169$-digit connection computation, the Stokes layout, and
the integer-relation battery, each with a canonical content hash recorded in
full in \texttt{HASH\_MANIFEST.json}; together with
the Lean project (buildable, with the axiom-audit file) and a snapshot of the
claims ledger. A \texttt{README} gives a one-command verification path per
artifact, and \texttt{verify.py} re-derives every embedded self-hash and file
hash against the manifest (exit code~$0$).

\medskip\noindent
The load-bearing frozen inputs are pinned here by SHA-256 prefix (full
$64$-hex digests in \texttt{HASH\_MANIFEST.json}). The $169$-digit
integer-relation null
\texttt{ebr3\_b\_pslq\_results.json} begins
\texttt{9a3f942def64737d{\dots}};
the connection coefficient
\texttt{cc4\_1\_connection\_results.json} begins
\texttt{f3400831cc964464{\dots}};
and the claims-ledger snapshot begins
\texttt{9e6f3fa5a2986cce{\dots}}.
(The ledger hash is non-circular: no claim references the bytes of this
document.)

\begin{thebibliography}{9}
\bibitem{EBRI} Papanokechi, \emph{An Edge--Borel Radius Theorem for Positive
Polynomial Continued Fractions}, Zenodo,
\href{https://doi.org/10.5281/zenodo.20564079}{10.5281/zenodo.20564079} (v1.0:
\texttt{20564080}).
\bibitem{EBRII} Papanokechi, \emph{The EBR Amplitude as a Connection
Coefficient: Characterization and a Rigidity Dividing-Line Conjecture},
Zenodo, \href{https://doi.org/10.5281/zenodo.20566465}{10.5281/zenodo.20566465}
(v1.2: \texttt{20571232}).
\bibitem{vdPutSinger} M.~van der Put and M.~F.~Singer, \emph{Galois Theory of
Linear Differential Equations}, Grundlehren der mathematischen Wissenschaften
\textbf{328}, Springer, 2003.
\bibitem{Andre} Y.~Andr\'e, \emph{$G$-functions and Geometry}, Aspects of
Mathematics \textbf{E13}, Vieweg, 1989.
\bibitem{Katz} N.~M.~Katz, \emph{Nilpotent connections and the monodromy
theorem: applications of a result of Turrittin}, Publ.\ Math.\ IH\'ES
\textbf{39} (1970), 175--232.
\bibitem{Kovacic} J.~J.~Kovacic, \emph{An algorithm for solving second order
linear homogeneous differential equations}, J.\ Symbolic Comput.\ \textbf{2}
(1986), 3--43.
\bibitem{PSLQ} H.~R.~P.~Ferguson, D.~H.~Bailey, and S.~Arno, \emph{Analysis of
PSLQ, an integer relation finding algorithm}, Math.\ Comp.\ \textbf{68} (1999),
351--369.
\bibitem{Mathlib} The Mathlib Community, \emph{The Lean Mathematical Library},
CPP 2020, 367--381.
\end{thebibliography}

\end{document}
